\documentclass[11pt]{article}
\usepackage[T1]{fontenc}
\usepackage{lmodern}
\usepackage[margin=1.05in]{geometry}
\usepackage{amsmath,amssymb,amsthm,mathtools}
\usepackage{microtype,booktabs,array}
\usepackage{xcolor}
\usepackage[colorlinks=true,linkcolor=blue!45!black,citecolor=blue!45!black,urlcolor=blue!45!black]{hyperref}
\usepackage[nameinlink,capitalise,noabbrev]{cleveref}
\usepackage{listings}
\lstset{basicstyle=\ttfamily\small,columns=fullflexible,keepspaces=true,breaklines=true,
  literate={≠}{{$\neq$}}1 {→}{{$\to$}}1 {∀}{{$\forall$}}1 {∃}{{$\exists$}}1 {¬}{{$\neg$}}1
           {∈}{{$\in$}}1 {∉}{{$\notin$}}1 {∧}{{$\wedge$}}1 {∨}{{$\vee$}}1 {⟨}{{$\langle$}}1
           {⟩}{{$\rangle$}}1 {×}{{$\times$}}1 {·}{{$\cdot$}}1 {•}{{$\bullet$}}1}
\numberwithin{equation}{section}
\newtheorem{theorem}{Theorem}[section]
\newtheorem{proposition}[theorem]{Proposition}
\newtheorem{lemma}[theorem]{Lemma}
\newtheorem{corollary}[theorem]{Corollary}
\theoremstyle{definition}
\newtheorem{definition}[theorem]{Definition}
\theoremstyle{remark}
\newtheorem{remark}[theorem]{Remark}
\newcommand{\F}{\mathbb F}
\newcommand{\N}{\mathbb N}
\DeclareMathOperator{\AGL}{AGL}
\title{New lower bounds for line-free sets in $\mathbb{F}_7^3$ and $\mathbb{F}_{11}^3$}
\author{Claude (AI model by Anthropic)\\ \small run and submitted by Aleksandr Bersenev}
\date{October 8, 2026\\[3pt]\footnotesize Produced autonomously by Claude (Opus 5.5) agents in Claude Code; run and submitted by Aleksandr Bersenev.\\ \footnotesize Not an official Anthropic publication. Not peer reviewed.}
\hypersetup{pdftitle={New lower bounds for line-free sets in F7^3 and F11^3},pdfauthor={Claude (AI model by Anthropic); submitted by Aleksandr Bersenev}}
\begin{document}
\maketitle
\begin{abstract}
A subset of $\F_p^n$ is \emph{line-free} if it contains no affine line, equivalently no
$p$-term arithmetic progression with nonzero common difference; $r_p(\F_p^n)$ denotes the
largest size of such a set. We exhibit an explicit line-free set of $227$ points in $\F_7^3$
and an explicit line-free set of $1010$ points in $\F_{11}^3$, so that
$r_7(\F_7^3)\ge 227$ and $r_{11}(\F_{11}^3)\ge 1010$; the previously best known lower bounds were
$225$ and $1005$. By the product construction, $r_7(\F_7^{3k})\ge 227^k$ for every $k\ge0$, and
hence $\lim_{n\to\infty} r_7(\F_7^n)^{1/n}\ge 227^{1/3}\approx 6.1002$; this improves the
three-dimensional product bound $225^{1/3}\approx 6.0822$ of Elsholtz et al., but not the best known
bound $7^{20/21}\approx 6.3805$, which follows from the line-evasive sets of F\"uhrer and Taranchuk
(2024). The sets were found by tabu search with iterated local search.
Both lower bounds and the statement $r_7(\F_7^{3k})\ge 227^k$ for all $k$ are certified in
Lean~4 (core library only) by kernel-checked computation and a structural proof; the only
step not formalised is the passage to the limit (Fekete's lemma).
\end{abstract}
\tableofcontents

\section{Introduction}

For a prime $p$ and integers $n\ge1$, $3\le k\le p$, let $r_k(\F_p^n)$ denote the maximum size of a
subset of $\F_p^n$ containing no $k$-term arithmetic progression $a,a+d,\dots,a+(k-1)d$ with
$d\neq0$. The case $k=p=3$ is the cap set problem. For $k=p$ a nontrivial $p$-term progression
is exactly an affine line $\{a+td:t\in\F_p\}$, so $r_p(\F_p^n)$ is the largest size of a
\emph{line-free} set. Equivalently, $p^n-r_p(\F_p^n)$ is the smallest size of a set of points of
the affine space $\mathrm{AG}(n,p)$ meeting every line (for $n=3$: a $2$-blocking set).

Line-free sets were studied systematically by Elsholtz, F\"uhrer, F\"uredi, Kov\'acs, Pach,
Simon and Velich~\cite{EFFKPSV}. The hypercube $\{0,1,\dots,p-2\}^n$ is line-free, so
$r_p(\F_p^n)\ge(p-1)^n$, with equality for $n\le 2$ (Jamison~\cite{Jamison}, Brouwer and
Schrijver~\cite{BrouwerSchrijver}, as quoted in~\cite{EFFKPSV}). In dimension three,
\cite{EFFKPSV} proves
\[
(p-1)^3+p-2\sqrt p\;\le\; r_p(\F_p^3)\;\le\; p^3-2p^2-(\sqrt2-1)p+2 \qquad(p\ge5),
\]
a further general bound $r_p(\F_p^n)\ge(p-1)^n+\frac{n-2}{2}(p-1)(p-2)^{n-3}$ (their Theorem~1.7,
which for $n=3$ gives $(p-1)^3+\frac{p-1}{2}$), the special bound $r_7(\F_7^3)\ge225$
(their Theorem~1.2 and the remark following its proof), a $70$-point line-free set in $\F_5^3$
found by a branch-and-cut computation, and the upper bounds $r_5(\F_5^3)<74$, $r_7(\F_7^3)<243$.
Kov\'acs~\cite{Kovacs} recently improved the lower bound to
$p^3-3p^2+\frac18p^{3/2}-\frac7{16}p=(p-1)^3+\frac18p^{3/2}-O(p)$; his construction uses
$\lfloor\sqrt p/4\rfloor$ special layers and is void for $p<16$, and the explicit bound exceeds
$(p-1)^3$ only for $p\ge757$, so it gives nothing for $p=7,11,13$. \cite{Kovacs} states that,
to the author's knowledge, there had been no further improvements of the bounds
of~\cite{EFFKPSV} from either side.

A further classical construction is the \emph{cone} (S.~Ball's lecture notes, as quoted in the
MathOverflow discussion~\cite{MO229442} opened by A.~Bishnoi in January 2016): if $B$ is a
double blocking set of $\mathrm{PG}(2,p)$, i.e.\ a set meeting every line in at least two points,
then the affine part of the union of the lines joining a fixed affine point to the points of $B$
meets every affine line, so that $r_p(\F_p^3)\ge p^3-1-(p-1)|B|$. With the sides of a triangle
($|B|=3p$) this is exactly $(p-1)^3$. Double blocking sets of size $3p-1$ in $\mathrm{PG}(2,p)$
are known for $p=13$ (Braun, Kohnert and Wassermann~\cite{BKW}; Csajb\'ok and
H\'eger~\cite{CsH}) and for $p=19,31,37,43$~\cite{CsH}; none exists for $p\le 8$ (Ball and
Blokhuis~\cite{BallBlokhuis}), and for $p=11$ none with fewer than $3p=33$ points is known
(\cite{CsH}; the lower bound is $31$~\cite{BallBlokhuis}). For $p=13$ the cone over the
$38$-point set of~\cite[\S4.1]{CsH} is therefore a line-free set of $13^3-1-12\cdot38=1740$
points (we verified this by computer; see \texttt{data/cone13\_1740.py}), which is larger than the
$1734$ points of
\cite[Thm.~1.1 and 1.7]{EFFKPSV}; \cref{tab:bounds} takes this into account. Zare's
single-point improvement $(p-1)^3+1$~\cite{MO229442}, mentioned in~\cite{EFFKPSV}, and the
constructions of affine $2$-blocking sets of Bishnoi, D'haeseleer, Gijswijt and
Potukuchi~\cite{BDGP}, which concern growing dimension, do not affect the small cases.

Since products of line-free sets are line-free (\cref{lem:product}), $r_p(\F_p^n)$ is
supermultiplicative in $n$, and the limit $\alpha_p=\lim_n r_p(\F_p^n)^{1/n}$ exists.
Small-dimensional examples give lower bounds for $\alpha_p$; \cite{EFFKPSV} obtained
$\alpha_5\ge 70^{1/3}\approx4.121$ and $\alpha_7\ge225^{1/3}\approx6.082$, which were better than
the bounds $4.090$ and $6.066$ (their Table~1) obtained from Frankl, Graham and R\"odl's
construction in dimension $2p$. These bounds have since been superseded: Theorem~1.1 of
F\"uhrer and Taranchuk~\cite{FuhrerTaranchuk} (first version March 2024) partitions
$\F_q^{n}$, $n=\binom{t+1}2$, into $q$ sets each meeting every line in at most $t$ points; for
$t=p-1$ these are line-free sets of size $p^{n-1}$ in $\F_p^{n}$, $n=p(p-1)/2$, whence
$\alpha_p\ge p^{1-2/(p^2-p)}$, i.e.\ $\alpha_5\ge5^{9/10}\approx4.2567$,
$\alpha_7\ge7^{20/21}\approx6.3805$ and $\alpha_{11}\ge11^{54/55}\approx10.5307$. Since
$4.2567^3>73$, $6.3805^3>242$ and $10.5307^3>1085$ exceed the upper bounds for
$r_p(\F_p^3)$, no three-dimensional set can reach these values, and the product bound
of \cref{cor:alpha} below is not a record.

\Cref{tab:bounds} summarises the state of the art for small $p$ before this note.

\begin{table}[h]
\centering
\begin{tabular}{@{}cccl@{}}
\toprule
$p$ & lower bound & upper bound & source of the lower bound\\
\midrule
5 & 70 & 73 & computer search (branch and cut) \cite{EFFKPSV}\\
7 & 225 & 242 & algebraic construction, \cite[Thm.~1.2]{EFFKPSV}\\
11 & 1005 & 1085 & $(p-1)^3+(p-1)/2$, \cite[Thm.~1.1 and 1.7]{EFFKPSV}\\
13 & 1740 & 1855 & cone over a $38$-point double blocking set of $\mathrm{PG}(2,13)$ \cite{BKW,CsH}\\
\bottomrule
\end{tabular}
\caption{Bounds for $r_p(\F_p^3)$ before this note. Upper bounds: \cite[Thm.~1.5, 1.6]{EFFKPSV}
for $p=5,7$ and the exact form of \cite[Cor.~1.4]{EFFKPSV} for $p=11,13$ (its simplified form
$p^3-2p^2-(\sqrt2-1)p+2$ gives $1086$ for $p=11$). The entry $1740$ for $p=13$ combines known
results (cone construction and \cite{BKW,CsH}) and is not stated explicitly in the literature we
found.}
\label{tab:bounds}
\end{table}

\paragraph{What is new.} We improve the lower bounds for $p=7$ and $p=11$
(\cref{thm:main7,thm:main11}) by two and five points respectively; the exact values remain open.
The product bound for $\alpha_7$ (\cref{cor:alpha}) is included for completeness only, since it
is weaker than the bound of~\cite{FuhrerTaranchuk}. To our knowledge, as of October~8, 2026,
$227$ and $1010$ are the best known lower bounds for $r_7(\F_7^3)$ and $r_{11}(\F_{11}^3)$; we
searched arXiv and the web for later work on line-free sets and on $2$-blocking sets in
$\mathrm{AG}(3,p)$ and found no improvement of the values in \cref{tab:bounds} (in particular
none in~\cite{Kovacs,FuhrerTaranchuk}; the cone construction gives $216$ for $p=7$ and at most
$1000$ for $p=11$ with the known double blocking sets). The sets were found by computer search; they are
irregular (the $227$-point set has trivial stabiliser in $\AGL(3,7)$) and we do not have a
conceptual description of them. The proofs consist of an exhaustive verification, which we
also certify in Lean~4.

\section{Definitions and results}

\begin{definition}
Let $p$ be a prime. A \emph{line} of $\F_p^n$ is a set $\{a+td: t\in\F_p\}$ with
$a,d\in\F_p^n$, $d\neq0$. A set $S\subseteq\F_p^n$ is \emph{line-free} if it contains no line.
$r_p(\F_p^n)$ is the maximum size of a line-free subset of $\F_p^n$, and
$\alpha_p=\lim_{n\to\infty}r_p(\F_p^n)^{1/n}$.
\end{definition}

\begin{theorem}\label{thm:main7}
The set $S_7\subseteq\F_7^3$ of $227$ points listed in \cref{app:set7} is line-free. Hence
$r_7(\F_7^3)\ge227$.
\end{theorem}

\begin{theorem}\label{thm:main11}
The set $S_{11}\subseteq\F_{11}^3$ of $1010$ points listed in \cref{app:set11} is line-free.
Hence $r_{11}(\F_{11}^3)\ge1010$.
\end{theorem}

\begin{lemma}[product construction]\label{lem:product}
If $A\subseteq\F_p^m$ and $B\subseteq\F_p^k$ are line-free, then
$A\times B\subseteq\F_p^{m+k}$ is line-free.
\end{lemma}

\begin{proof}
Let $L=\{a+td\}$ be a line of $\F_p^{m+k}$ and write $a=(a_1,a_2)$, $d=(d_1,d_2)$. If $d_1\neq0$,
the projection $\{a_1+td_1\}$ is a line of $\F_p^m$, so some $a_1+t_0d_1\notin A$, and then
$a+t_0d\notin A\times B$. If $d_1=0$ then $d_2\neq0$ and the same argument applies to the second
factor.
\end{proof}

\begin{corollary}\label{cor:power}
For every $k\ge0$, $r_7(\F_7^{3k})\ge227^k$ and $r_{11}(\F_{11}^{3k})\ge1010^k$.
\end{corollary}

\begin{proof}
Apply \cref{lem:product} $k-1$ times to $S_7$ (resp.\ $S_{11}$); $|S^k|=|S|^k$.
\end{proof}

\begin{corollary}\label{cor:alpha}
$\alpha_7\ge227^{1/3}\approx 6.1002$.
\end{corollary}

\begin{proof}
By \cref{lem:product}, $r_p(\F_p^{m+k})\ge r_p(\F_p^m)\,r_p(\F_p^k)$, so $\log r_p(\F_p^n)$ is
superadditive and by Fekete's lemma $\alpha_p=\lim_n r_p(\F_p^n)^{1/n}=\sup_n
r_p(\F_p^n)^{1/n}$. Taking $n=3$ and \cref{thm:main7} gives the claim.
\end{proof}

\begin{remark}
\Cref{cor:alpha} is not the best known bound: \cite[Thm.~1.1]{FuhrerTaranchuk} gives
$\alpha_7\ge7^{20/21}\approx6.3805$ and $\alpha_{11}\ge11^{54/55}\approx10.5307$ (see the
introduction), and no three-dimensional set can reach these values because $6.3805^3>242$ and
$10.5307^3>1085$. For $p=11$ the bound $1010^{1/3}\approx10.0332$ does not even reach the older
bound $\alpha_{11}\ge 11^{1/22}10^{21/22}\approx10.0434$ of Frankl, Graham and R\"odl
(\cite[Table~1]{EFFKPSV}), for which a line-free set of at least $1014$ points in $\F_{11}^3$
would be needed.
\end{remark}

\section{The two sets}

\subsection{Encoding and verification}
A point $(x,y,z)\in\F_p^3$ is encoded by the integer $x+py+p^2z\in[0,p^3)$, with coordinates in
$\{0,\dots,p-1\}$. The data files \texttt{data/linefree\_F7\_227.txt} and
\texttt{data/linefree\_F11\_1010.txt} list the encoded points (the \texttt{\_xyz} files list
coordinates); SHA-256 hashes are in \texttt{data/SHA256SUMS}:
\begin{center}\small
\begin{tabular}{@{}ll@{}}
\texttt{linefree\_F7\_227.txt} & \texttt{fc505c9d673a4d390ba665a9333d35998ea1bfa34ab03e9b2c6cbad9d5623636}\\
\texttt{linefree\_F11\_1010.txt} & \texttt{3c65acbb753817d20bc5a7fc118cddef9d43ead14d0ee5be425febdc11070010}
\end{tabular}
\end{center}
The sets are also drawn layer by layer in \cref{app:set7,app:set11}.

\begin{proof}[Proof of \cref{thm:main7,thm:main11}]
$\F_p^3$ has $p^2(p^2+p+1)$ lines: $2793$ for $p=7$ and $16093$ for $p=11$. For each base point
$a$ and each nonzero direction $d$ one checks that some $a+td$ lies outside the set. This finite
verification was carried out by three independent programs: the Python script
\texttt{data/verify.py} (all pairs $(a,d)$, $d\ne0$), a C implementation inside the search
program, and the Lean kernel (\cref{sec:lean}); an independent script written by a second party
also confirmed that no line is contained in either set.
\end{proof}

\subsection{Statistics}
Let $n_j$ be the number of lines meeting the set in exactly $j$ points. Then
\begin{align*}
S_7:&\quad (n_0,\dots,n_7)=(26,34,112,282,603,993,743,\mathbf 0),\\
S_{11}:&\quad (n_0,\dots,n_{11})=(91,99,58,19,18,103,503,1954,4898,5909,2441,\mathbf 0).
\end{align*}
(As a consistency check, $\sum_j jn_j=|S|(p^2+p+1)$, e.g.\ $12939=227\cdot57$.)

Both sets are maximal with respect to inclusion: adding any of the missing points creates a full
line. Concerning structure:
\begin{itemize}
\item In the parallel class of planes $x+2y+3z=c$ of $\F_7^3$, the planes meet $S_7$ in
$14,34,35,36,36,36,36$ points ($36=(p-1)^2=r_7(\F_7^2)$ is the maximum for a plane). The
algebraic $225$-point set of~\cite{EFFKPSV} consists of six layers of $36$ points and one layer
of $9$ points. $105$ points of $S_7$ have a ``local line profile'' (the multiset of
$|L\cap S_7|$ over the $57$ lines $L$ through the point) shared by no other point of $S_7$; these
points span $\F_7^3$ affinely, so the stabiliser of $S_7$ in $\AGL(3,7)$ is trivial.
\item In the parallel class $x+7y+7z=c$ of $\F_{11}^3$, ten planes contain $100=(p-1)^2$ points
of $S_{11}$ and one plane contains $10$, i.e.\ $S_{11}$ has the coarse shape ``ten maximal planes
plus a sparse plane'' of hypercube-type constructions, with $10$ extra points instead of the
$(p-1)/2=5$ of~\cite[Thm.~1.7]{EFFKPSV}.
\end{itemize}

\section{How the sets were found}\label{sec:search}

We searched for $k$-subsets of $\F_p^3$ minimising the number of fully contained lines.
\begin{enumerate}
\item \emph{Tabu search.} A swap move removes a point of the current set and adds a point of
the complement. Gains are maintained incrementally (for a point of the set, the number of full
lines through it; for a point outside, the number of lines through it with $p-1$ points of the
set); the best swap is chosen among all pairs, with random tie-breaking, short random tabu
tenures and aspiration. Many \emph{short} restarts work best: at $p=7$, $k=225$, restarts of
$400$--$800$ iterations gave about $40$ solutions per $150$ CPU-seconds, against $12$ with
$3000$ and $10$ with $6000$ iterations; a two-stage move rule, adaptive line weights and greedy
initialisation were all worse. With this method $224$ points are found in about one restart
in eight, $225$ in one in two hundred, and a $226$-point set was found after about five CPU-minutes;
no $227$-point set appeared in a further hour. For $p=11$, sets of size $1006$--$1008$ were found in
one to five minutes.
\item \emph{Iterated local search.} Restart the same tabu search from a known solution with
$10$--$30$ randomly chosen points exchanged and the target size increased by one. This is much
stronger: from a $225$-point set, $226$-point sets were found $61$ times in $150$ CPU-seconds
(with $10$ exchanged points), against once for random restarts. From the $226$-point set,
$S_7$ was found after about two CPU-minutes ($20$ exchanged points; $S_7$ shares $199$ points
with the $226$-point set). For $p=11$, chaining this step from a $1008$-point set gave
$1009$ and then $S_{11}$.
\end{enumerate}
A SAT-based neighbourhood test (re-optimising exactly all points in one plane, or in two parallel
planes, with the CaDiCaL solver) showed that the $226$-point set cannot be enlarged by such local
changes, which is consistent with the need for larger perturbations.

\paragraph{Negative results (not theorems).}
These are reported only to document the search effort.
\begin{itemize}
\item $p=7$: no $228$-point set in about $85$ CPU-minutes of iterated local search, including a
random walk over $227$-point solutions.
\item $p=11$: no $1011$-point set in about $47$ CPU-minutes.
\item $p=5$: $70$-point sets are found within a few hundred iterations, but no $71$-point set was
found (several CPU-minutes of tabu search and iterated local search). We do not know whether
$r_5(\F_5^3)=70$.
\item $p=13$: sets of size $1731$ were found within two minutes, but none of size $1735$ (the
bound $1734$ of~\cite{EFFKPSV} plus one) in about $35$ CPU-minutes of random restarts; iterated
local search was not tried for $p=13$. Note that the cone construction gives $1740$ points
(see the introduction), so the search did not even reach the best known value for $p=13$.
\item In $\F_5^4$ the search reached only $280=70\cdot4$ points, so it gave no improvement of
$\alpha_5$.
\end{itemize}

\section{Formal verification}\label{sec:lean}

Both theorems are proved in Lean~4.34.1 using only the core library and \texttt{Std} (no
Mathlib), in the self-contained files \texttt{lean/Main.lean} ($p=7$) and
\texttt{lean/Main11.lean} ($p=11$). No \texttt{sorry}, \texttt{axiom} or \texttt{native\_decide}
is used. The main statements and all definitions they depend on (for $p=7$; the $p=11$ file is
identical with $11$ and $1010$) are:
\begin{lstlisting}
abbrev Pt := Fin 7 × Fin 7 × Fin 7
def add (a b : Pt) : Pt := (a.1 + b.1, a.2.1 + b.2.1, a.2.2 + b.2.2)
def smul (t : Fin 7) (d : Pt) : Pt := (t * d.1, t * d.2.1, t * d.2.2)
def LineFree (L : List Pt) : Prop :=
  ∀ a d : Pt, d ≠ (0, 0, 0) → ∃ t : Fin 7, add a (smul t d) ∉ L
theorem main : ∃ L : List Pt, L.Nodup ∧ L.length = 227 ∧ LineFree L

def vadd (a b : List (Fin 7)) : List (Fin 7) := List.zipWith (· + ·) a b
def vsmul (t : Fin 7) (d : List (Fin 7)) : List (Fin 7) := d.map (t * ·)
def LineFreeN (n : Nat) (L : List (List (Fin 7))) : Prop :=
  ∀ a d : List (Fin 7), a.length = n → d.length = n → (∃ x ∈ d, x ≠ 0) →
    ∃ t : Fin 7, vadd a (vsmul t d) ∉ L
theorem main_pow (k : Nat) : ∃ M : List (List (Fin 7)), M.Nodup ∧ M.length = 227 ^ k ∧
    (∀ v ∈ M, v.length = 3 * k) ∧ LineFreeN (3 * k) M
\end{lstlisting}

\paragraph{Proof architecture.} The witness \texttt{theSet} is the explicit list of points,
sorted lexicographically. A bitmask with bit $p^2x+py+z$ set for each point is tied to it by
the kernel-checked equation \texttt{theSet = allPts.filter inMask}, where \texttt{allPts}
lists all $p^3$ points. \texttt{allPts} has no duplicates, because it maps onto
\texttt{List.range}~$p^3$ under the lexicographic index; this gives \texttt{theSet.Nodup}.
The length is checked by \texttt{decide +kernel}. For line-freeness, a Boolean checker
\texttt{checkAt a} is evaluated by \texttt{decide +kernel} on chunks of base points (7 slices
for $p=7$, 121 rows for $p=11$); the chunks are glued by \texttt{match} on \texttt{Fin} and
combined with soundness lemmas (\texttt{List.all\_eq\_true}, \texttt{List.any\_eq\_true},
membership of every point in \texttt{allPts}). For $p=7$ the checker tries all $342$ nonzero
directions. For $p=11$ it tries only directions whose first nonzero coordinate is $1$. The
reduction is proved in the file: a kernel check shows that $u^{-1}d$ is normalised for every
$d\ne0$, where $u$ is the leading coordinate of $d$, and associativity of multiplication in
\texttt{Fin 11} (checked by \texttt{decide +kernel}) gives $s(u^{-1}d)=(su^{-1})d$.
\texttt{main\_pow} is a structural proof: the product of two line-free, duplicate-free lists of
vectors of fixed lengths is line-free and duplicate-free, of the product size
(\cref{lem:product}). It is proved by induction on $k$ and instantiated with \texttt{main}'s
witness.

\paragraph{Statement audit.}
(i) \texttt{Fin 7} addition and multiplication are taken modulo $7$, so \texttt{Pt} with
\texttt{add} and \texttt{smul} is the vector space $\F_7^3$.
(ii) \texttt{LineFree L} says: for every base point $a$ and every direction $d\neq0$, some
point $a+td$ is not in \texttt{L}, i.e.\ \texttt{L} contains no line. All pairs $(a,d)$ are
quantified; the direction normalisation for $p=11$ happens only inside the proof.
(iii) \texttt{L.Nodup} together with \texttt{L.length = 227} says that \texttt{L} consists of
$227$ distinct points.
(iv) In \texttt{main\_pow}, vectors of $\F_7^{3k}$ are lists of length $3k$, and this length is
part of the statement; ``$d\neq0$'' is \texttt{$\exists$ x $\in$ d, x $\neq$ 0}. For $k=0$ the
statement is the trivial $r_7(\F_7^0)\ge1$.
(v) Both theorems are existential statements whose witness is explicit, so there are no
hypotheses that could be vacuous. A separate control file
(\texttt{lean/NegTest.lean}, \texttt{\#eval} only) checks that the same checker rejects the set
with the line $\{(t,0,0)\}$ added, and the set with any single missing point added.

\paragraph{Axioms and resources.} \texttt{\#print axioms} reports, for both files,
\begin{lstlisting}
'main' depends on axioms: [propext, Classical.choice, Quot.sound]
'main_pow' depends on axioms: [propext, Classical.choice, Quot.sound]
\end{lstlisting}
These are the standard axioms; no \texttt{Lean.ofReduceBool} appears. Checking takes about
$140$\,s with peak memory $2.2$\,GB for $p=7$, and about $550$\,s with $0.96$\,GB for $p=11$,
on a shared machine. The outputs are in \texttt{lean/lean\_output.txt} and
\texttt{lean/lean\_output11.txt}.

\section{Limitations and open questions}
\begin{itemize}
\item Novelty is claimed to our knowledge as of October~8, 2026, based on a literature search;
unpublished improvements cannot be excluded.
\item The exact values remain open: $227\le r_7(\F_7^3)\le242$ and
$1010\le r_{11}(\F_{11}^3)\le1085$. The search gives no evidence that our sets are optimal.
\item The bound $\alpha_7\ge227^{1/3}$ of \cref{cor:alpha} is weaker than the known bound
$\alpha_7\ge7^{20/21}$ of~\cite{FuhrerTaranchuk}; for $p=5,7,11$ no three-dimensional
line-free set can improve the bounds for $\alpha_p$ that follow from~\cite{FuhrerTaranchuk}.
\item Not formalised: the existence of the limit $\alpha_7$ (Fekete's lemma), and the comparison
with the literature.
\item The sets are unstructured. A structural explanation of why random search beats the
algebraic constructions for $p=7,11$ might lead to improved general constructions of the form
$(p-1)^3+cp^{\beta}$. For $p=11$, about ten points beyond the hypercube fit into one sparse
plane.
\item Is $r_5(\F_5^3)=70$? This seems within reach of a careful exhaustive search with symmetry
reduction, but a formal certificate would be substantially harder.
\end{itemize}

\appendix
\section{The 227-point line-free set in \texorpdfstring{$\F_7^3$}{F7^3}}\label{app:set7}
Layers $z=0,\dots,6$; in each layer the rows are $y=0,\dots,6$ (top to bottom) and the columns
are $x=0,\dots,6$ (left to right). The symbol \texttt{\#} marks a point of $S_7$; the layer sizes
are given in parentheses.
{\small
\begin{verbatim}
z=0      z=1      z=2      z=3      z=4      z=5      z=6    
..#.###  #.###.#  #..####  #.#.###  #..#.##  #.#.#..  #.##.#.
####.#.  ###.##.  .#.####  #..####  .######  #.#.#..  ##.##.#
####..#  ####.##  .##.###  #....#.  ###.##.  ###..##  ####.##
##.###.  ##.#...  ##.#.##  .######  ##.####  ....##.  #.##.##
.####.#  ####.#.  ..###.#  #..####  ##.####  ##..###  ..##.##
#.###.#  ##.###.  ..####.  ###.#.#  .#.#..#  ###....  #.#####
##..###  ..#.###  .###.#.  #.#####  #..##.#  .#..###  ###.##.
(34)     (33)     (32)     (34)     (34)     (25)     (35)   
\end{verbatim}
}

\section{The 1010-point line-free set in \texorpdfstring{$\F_{11}^3$}{F11^3}}\label{app:set11}
Same conventions, with layers $z=0,\dots,10$, rows $y=0,\dots,10$ and columns $x=0,\dots,10$.
{\footnotesize
\begin{verbatim}
z=0          z=1          z=2          z=3          z=4          z=5        
###.##.#.##  #.#######.#  #####.##...  #########.#  ##.###..#..  .##.##.####
#..#.######  ##.##.#####  ######.##.#  ##.####.###  ##.###.####  #########..
##.##.####.  ###.#####.#  ##.#.######  #.##...###.  ######..##.  ###.###.###
..######..#  ##.########  ##.###.####  ###.######.  ###..######  ##.##.#.###
##.######.#  #####..####  .#########.  #.#.###.###  ##.####.##.  .#..#######
#######..##  #######..#.  #.#.###.#.#  ##.####.##.  .##########  #.#..####.#
#####.####.  ######..###  ######..###  .#######.##  #########.#  ########.#.
###.###.##.  ####.##.#.#  .###..####.  ####...####  ##.####...#  ..######.##
##.####.###  .#.#####..#  ###..######  #.##.###.#.  #.#.#.#####  .####..####
.#########.  .###.####..  #.######.##  #..#.#.####  #.#.#.####.  ####..##..#
####.##..##  ########..#  #.#.######.  .###....##.  #.###..##.#  ##.########
(92)         (94)         (92)         (86)         (89)         (91)       
\end{verbatim}
\begin{verbatim}
z=6          z=7          z=8          z=9          z=10       
#.##..####.  ##.#######.  ###.#######  #######.###  ####..#####
###.######.  ..##.######  ..#########  ##...######  ###.####.#.
####.##..##  .########.#  ####.######  .##########  ...########
######.####  ##...##.###  ########.##  #.#..##.#..  #####.##.#.
##.#.######  #####.##.##  #..###.#.##  #####.#.###  ######.##.#
########.##  #.#.#######  #####.##.##  #.###.#.#.#  ##.###.####
..#.#######  #.###..####  ##.#.##....  #...##.####  .###.#..###
#####.###.#  #####.#.#..  .#.########  #.####.####  ##.#######.
#.##.##.#.#  ##.#####.##  ######.##.#  #..#.#####.  .##..##.###
#####.#####  #.####.####  #######.##.  ###.#.##.##  #######.###
###.##.####  ########.#.  ###.#.#.###  ######..###  .######.###
(96)         (93)         (95)         (89)         (93)       
\end{verbatim}
}

\begin{thebibliography}{99}
\bibitem{EFFKPSV}
C.~Elsholtz, J.~F\"uhrer, E.~F\"uredi, B.~Kov\'acs, P.~P.~Pach, D.~G.~Simon, N.~Velich,
\emph{Maximal line-free sets in $\mathbb F_p^n$}, Period.\ Math.\ Hungar.\ \textbf{90} (2025),
7--21, \url{https://doi.org/10.1007/s10998-024-00617-x}; arXiv:2310.03382 (v2, May 2024).
\url{https://arxiv.org/abs/2310.03382}

\bibitem{Kovacs}
B.~Kov\'acs, \emph{A superlinear improvement on line-free sets in $\mathbb F_p^3$},
arXiv:2605.23437 (May 2026, revised September 2026). \url{https://arxiv.org/abs/2605.23437}

\bibitem{FuhrerTaranchuk}
J.~F\"uhrer, V.~Taranchuk, \emph{Large line-free sets and their applications},
arXiv:2403.18611 (v1 March 2024, v2 February 2026). \url{https://arxiv.org/abs/2403.18611}

\bibitem{MO229442}
A.~Bishnoi (question), D.~Zare (answer), \emph{Blocking sets in three dimensional finite affine
spaces}, MathOverflow, January 2016. \url{https://mathoverflow.net/q/229442}

\bibitem{BallBlokhuis}
S.~Ball, A.~Blokhuis, \emph{On the size of a double blocking set in $PG(2,q)$},
Finite Fields Appl.\ \textbf{2} (1996), 125--137.

\bibitem{BKW}
M.~Braun, A.~Kohnert, A.~Wassermann, \emph{Construction of $(n,r)$-arcs in $PG(2,q)$},
Innov.\ Incidence Geom.\ \textbf{1} (2005), 133--141.

\bibitem{CsH}
B.~Csajb\'ok, T.~H\'eger, \emph{Double blocking sets of size $3q-1$ in $PG(2,q)$},
European J.\ Combin.\ \textbf{78} (2019), 73--89. arXiv:1805.01267.

\bibitem{BDGP}
A.~Bishnoi, J.~D'haeseleer, D.~Gijswijt, A.~Potukuchi, \emph{Blocking sets, minimal codes and
trifferent codes}, arXiv:2301.09457 (2023).

\bibitem{Jamison}
R.~E.~Jamison, \emph{Covering finite fields with cosets of subspaces},
J.\ Combin.\ Theory Ser.\ A \textbf{22} (1977), 253--266.

\bibitem{BrouwerSchrijver}
A.~E.~Brouwer, A.~Schrijver, \emph{The blocking number of an affine space},
J.\ Combin.\ Theory Ser.\ A \textbf{24} (1978), 251--253.
\end{thebibliography}

\end{document}
